\chapter{Five Source Families and the Archimedean Tate Positivity Residual}
\label{v4-arith:w8-atp:chapter}

\emph{The reduction of Chapter~\ref{v4-arith:w5-synth:chapter} places
the Riemann hypothesis in the form of archimedean Tate positivity,
A-TP (\Cref{v4-arith:w5-a:conj:archimedean-tate-positivity}). Five
source families meet this same Weil inequality from different sides:
Arthur--Selberg scattering, geometric Langlands, Iwasawa--Kato,
arithmetic quantum chaos, and Arakelov--Faltings--Hriljac. None of
them removes the Weil positivity problem. Pair-correlation controls
statistics after the zero-line input; it does not supply the zero-line
input. Iwasawa theory controls finite-prime symmetry; it does not
supply the archimedean sign. Arakelov height positivity gives useful
sub-class inequalities only after the Bott--Chern and vertical terms
are controlled.}

\section{A-TP: Statement and Independent Obstruction}
\label{v4-arith:w8-atp:target}

Let \(\mathrm{PW}(\mathbb{R})\) be the Paley-Wiener
space of even Schwartz functions extending holomorphically to a strip
with \((1+|z|)^{-2-\delta}\) decay; let \(\tilde f(x)=\overline{f(-x)}\).
The Weil distribution \(W\) on \(\mathrm{PW}(\mathbb{R})\) is the linear
functional of \Cref{v4-arith:tfe:thm:weil}:
\begin{equation}
\label{v4-arith:w8-atp:eq:weil-W}
W(h)=
\widehat h\bigl(\tfrac{1}{2i}\bigr)+
\widehat h\bigl(-\tfrac{1}{2i}\bigr)
-W_\infty(h)
-2\sum_{p}\sum_{k\geq 1}\frac{\log p}{p^{k/2}}h(k\log p).
\end{equation}
A-TP asserts \(W(f*\tilde f)\geq 0\) for all \(f\in\mathrm{PW}(\mathbb{R})\),
equivalent to RH via Weil 1952 (\emph{Bull.~Soc.~Math.~France} 80).

The five source lists below are pairwise disjoint from each other
and from those of
Chapters~\ref{v4-arith:w5-a:chapter}--\ref{v4-arith:w7-j:ATP-adelic-scattering},
modulo the common ground of Tate's thesis and Weil's explicit formula:

\begin{itemize}
\item \textbf{Arthur-Selberg truncation}: Arthur 1978--2013 trace-formula
series; Jacquet-Langlands 1970; Selberg 1956;
Ng\^{o} 2010 \emph{Publ.~IHES} 111;
Flicker-Kazhdan; M\oe{}glin-Waldspurger.
\item \textbf{Geometric Langlands and Fargues--Scholze}:
Beilinson--Drinfeld 2004; Arinkin-Gaitsgory arXiv:1201.6343;
Fargues--Scholze arXiv:2102.13459; Lafforgue arXiv:1209.5352; Drinfeld
1983, 1987; Deligne \emph{Weil II}; Frenkel-Gaitsgory.
\item \textbf{Iwasawa-Kato}: Mazur-Wiles 1984; Kato 2004
\emph{Ast\'{e}risque} 295; Perrin-Riou 1995; Fukaya-Kato 2006; Colmez;
Bloch-Kato 1990; Kolyvagin 1988; Coates-Wiles 1977; Ferrero-Washington 1979.
\item \textbf{Arithmetic quantum chaos}: Bohigas-Giannoni-Schmit
1984; Berry-Keating 1999; Montgomery 1973; Odlyzko 1987; Rudnick-Sarnak
1996; Katz-Sarnak 1999; Keating-Snaith 2000; Lindenstrauss 2006;
Soundararajan 2010; Sarnak 1995; Hejhal 1994.
\item \textbf{Arakelov--Faltings-Hriljac}: Arakelov 1974;
Faltings 1984; Hriljac 1985; Gillet-Soul\'{e} 1990; Bost 1999;
Bost-K\"{u}nnemann 2010; Bombieri 1971; Gross-Zagier 1986; Beilinson 1985;
Vojta 1987; Bismut-Gillet-Soul\'{e} 1988.
\end{itemize}

Each list's intersection with those of
Chapters~\ref{v4-arith:w5-a:chapter}--\ref{v4-arith:w7-j:ATP-adelic-scattering}
is empty modulo Tate's thesis and Weil's explicit formula.

\section{Five Source Families}
\label{v4-arith:w8-atp:convergence}

\begin{theorem}[five source-family boundary for A-TP]
\label{v4-arith:w8-atp:thm:convergence}
\ClaimStatusProvedHere. The five source families below do not give an
independent proof of A-TP. Each lands on one of three mathematical
forms: the Weil inequality itself, a conditional sub-class
diagnostic, or an auxiliary structure which has no archimedean
positivity content without an additional RH-strength input.
\end{theorem}

\begin{center}
\small
\begin{tabular}{lp{7cm}}
\toprule
Tradition & Exact mathematical form \\
\midrule
Arthur--Selberg & Maass--Selberg scattering recovers the
archimedean gamma logarithmic derivative inside \(W_\infty\). It does
not remove the finite-prime and polar terms of \(W\). \\
Geometric Langlands & The spectral-action language can encode
tempered parameters after the local and global comparison inputs. It
does not prove Weil positivity for the zeta divisor. \\
Iwasawa--Kato & Euler systems and main-conjecture structures give
finite-prime arithmetic control. They carry no theorem forcing the
archimedean Weil sign. \\
Arithmetic quantum chaos & GUE and pair-correlation statements are
statistics of ordinates after a zero-line model is fixed. They do not
replace \(z_\rho=-i(\rho-\tfrac12)\) by a real ordinate. \\
Arakelov--Faltings--Hriljac &
The arithmetic Hodge index gives negative definiteness on
degree-zero height pairings. Weil positivity follows only after the
Bott--Chern and vertical residual terms have the required sign. \\
\bottomrule
\end{tabular}
\end{center}

\begin{definition}[consolidated residual]
\label{v4-arith:w8-atp:defn:consolidated-residual}
\ClaimStatusDefinitional. The \emph{A-TP residual} is the full Weil
quadratic inequality
\[
W(f*\widetilde f)\geq 0,\qquad f\in\mathrm{PW}(\mathbb R),
\]
with polar, finite-prime, and archimedean terms all present. The
archimedean digamma kernel is one summand of this residual, not a
replacement for it. Denote the full residual by
\(\mathsf{R}_{\mathrm{A\text{-}TP}}\).
\end{definition}

\begin{remark}[one obstruction, several shadows]
\label{v4-arith:w8-atp:rem:fivefold}
The source families give shadows of the same Weil inequality. A
shadow is not an equivalent coordinate chart unless it retains the
polar terms, the von Mangoldt sum, the archimedean gamma term, and
the full Paley--Wiener cone.
\end{remark}

\section{Two Conditional A-TP Diagnostics}
\label{v4-arith:w8-atp:new-closures}

\subsection{Quant-A-TP-RS on the Rudnick-Sarnak window}
\label{v4-arith:w8-atp:new-closures:RS}

\begin{theorem}[Rudnick--Sarnak window does not give unconditional A-TP]
\label{v4-arith:w8-atp:thm:Quant-A-TP-RS}
\ClaimStatusNeedsVerification. Let \(\mathrm{PW}^{\mathrm{CP}}\subset\mathrm{PW}(\mathbb{R})\)
be the co-Poisson sub-class (functions whose translates are
orthogonal to each co-Poisson form; see \Cref{v4-arith:tfe:thm:equiv}).
On the restricted class
\begin{equation}
\label{v4-arith:w8-atp:eq:PW-CP-RS}
\mathrm{PW}^{\mathrm{CP}}\cap\bigl\{f:\operatorname{supp}\hat f\subset[-1,1]\bigr\}
\end{equation}
(the co-Poisson class meeting a small spectral window), there is no
unconditional proof of
\(W(f*\tilde f)\geq0\). Pair-correlation estimates control averaged
spacing statistics; they do not turn the Weil zero functional into the
positive square sum \(\sum_\rho|\hat f(\gamma_\rho)|^2\) without RH.
\end{theorem}

\begin{proof}[Sketch]
Unconditionally the zero side is evaluated on the divisor
\(z_\rho=-i(\rho-\tfrac12)\), not on real ordinates
\(\gamma_\rho\). The replacement by a positive square sum is exactly
the RH-level specialization. Rudnick--Sarnak pair correlation does not
remove that specialization.
\end{proof}

\begin{remark}[domain of the Rudnick-Sarnak sub-closure]
\label{v4-arith:w8-atp:rem:Quant-RS-narrow}
The sub-class \(\mathrm{PW}^{\mathrm{CP}}\cap\{\operatorname{supp}\hat f
\subset[-1,1]\}\) has measure zero in \(\mathrm{PW}(\mathbb{R})\)
under any natural Gaussian measure: the Rudnick-Sarnak window
\(\{|t|\leq 1\}\) is strictly smaller than any dense sub-class.
\Cref{v4-arith:w8-atp:thm:Quant-A-TP-RS} is disjoint from the
sub-closures of
Chapters~\ref{v4-arith:w5-a:chapter}--\ref{v4-arith:w7-j:ATP-adelic-scattering},
which give sub-class inequalities parameterized by central-value and
bounded-kernel methods.
\end{remark}

\subsection{Arakelov-Heegner diagnostic on \(\mathrm{PW}^{\mathrm{Heeg}}_F\)}
\label{v4-arith:w8-atp:new-closures:heeg}

\begin{theorem}[Arakelov-Heegner A-TP diagnostic]
\label{v4-arith:w8-atp:thm:Arakelov-Heegner}
\ClaimStatusProvedHereConditional \emph{(on Faltings-Hriljac 1984/1985
arithmetic Hodge index, Gross-Zagier 1986 central-derivative formula,
Jorgenson-Kramer 2003 Petersson-vs-Arakelov Bott-Chern comparison,
and the sign of the residual term \(R_{F,K}(f)\))}.
Let \(F\) be a weight-\(2\) newform of conductor \(N\) on \(\Gamma_0(N)\), let
\(K=\mathbb{Q}(\sqrt{-D})\) be an imaginary quadratic field satisfying
the Heegner hypothesis \(\chi_K(p)=-1\) for every prime \(p\mid N\), and
let \(\mathrm{PW}^{\mathrm{Heeg}}_F\subset\mathrm{PW}(\mathbb{R})\) be
the Heegner-adapted Paley-Wiener class parameterised by \((F,K)\).
Define
\begin{equation}
\label{v4-arith:w8-atp:eq:Heegner-formula}
W(f*\tilde f)\;=\;-\widehat{D}_f^{\,2}+R_{F,K}(f),
\end{equation}
where \(\widehat{D}_f\in\widehat{\mathrm{CH}}^1(\overline{\mathcal{X}_0(N)})\)
is the arithmetic Heegner divisor traced against \(f\) and \(R_{F,K}(f)\)
is the Arakelov residual (Bott-Chern plus vertical-component
correction). If the degree-zero projection satisfies
\(\widehat{D}_f^{\,2}\le0\) and \(R_{F,K}(f)\ge0\), then
\(W(f*\tilde f)\geq0\) on that Heegner-adapted sub-class. The theorem
is therefore a diagnostic, not an unconditional A-TP theorem.
\end{theorem}

\begin{proof}[Sketch]
Faltings-Hriljac 1984/1985 arithmetic Hodge index on
\(\widehat{\mathrm{CH}}^1(\overline{\mathcal{X}_0(N)})\) has signature
\((1,+\infty)\), with the negative-definite part on degree-zero classes
modulo vertical components (N\'{e}ron-Tate). Gross-Zagier 1986 expresses
\(L'(1,F)\) as the N\'{e}ron-Tate height \(\widehat{h}(y_K)\) of a Heegner
point \(y_K\in X_0(N)(K)\); N\'{e}ron-Tate heights are non-negative.
The sign of the Bott-Chern and vertical residual is not automatic for
an arbitrary Paley-Wiener trace. It is precisely the extra hypothesis
\(R_{F,K}(f)\ge0\). Under these hypotheses the displayed identity gives
Weil positivity.
\end{proof}

\begin{remark}[disjointness from the Rankin-Selberg diagnostic]
\label{v4-arith:w8-atp:rem:heeg-distinct}
\Cref{v4-arith:kp:thm:main-full} of Chapter~\ref{v4-arith:kp-compat}
provides the Rankin-Selberg central-value diagnostic through
\(\Upsilon_F\) at \(L(1/2,F\otimes\bar F)\).
\Cref{v4-arith:w8-atp:thm:Arakelov-Heegner} gives a non-overlapping
Heegner diagnostic through Gross-Zagier's \emph{central derivative}
\(L'(1,F)\). The two diagnostics address disjoint analytic sectors of
the same newform \(F\); neither removes the full A-TP residual without
its stated sign hypotheses.
\end{remark}

\section{Two Negative Results: Iwasawa Auxiliary Structure and BGS Implication Direction}
\label{v4-arith:w8-atp:negative}

\subsection{Iwasawa-Kato Orthogonality to A-TP}
\label{v4-arith:w8-atp:negative:iwasawa}

\begin{theorem}[Iwasawa-Kato orthogonality to A-TP]
\label{v4-arith:w8-atp:thm:iwasawa-auxiliary-structure}
\ClaimStatusProvedHereConditional \emph{(on completeness of the
Iwasawa-Kato list: Mazur-Wiles 1984, Kato 2004 Ast\'{e}risque 295,
Perrin-Riou 1995, Fukaya-Kato 2006, Colmez, Bloch-Kato 1990,
Kolyvagin 1988, Coates-Wiles 1977, Ferrero-Washington 1979,
Connes-Marcolli, Skinner-Urban)}.
No method in the standard Iwasawa-Kato list closes A-TP.
Specifically: Mazur-Wiles value-positivity, Kato Euler system plus
Bloch-Kato-Beilinson height pairing, Ferrero-Washington archimedean
analog, Perrin-Riou plus Bloch-Kato exponential, the \(p\to\infty\)
Connes-Marcolli limit, Kato reciprocity maps \(\exp^{*}/\log\),
Kolyvagin plus Coates-Wiles height positivity, Colmez-Fontaine
\(p\)-adic Langlands, Fukaya-Kato non-commutative plus Skinner-Urban,
and concrete-Dirichlet \(p\)-adic transport each fail to deliver
A-TP positivity beyond sectors already closed by other means.
Iwasawa theory supplies the functional-equation symmetry
\(s\leftrightarrow 1-s\) at every finite prime
(as in \Cref{v4-arith:thm:P1-resolution}) but carries no positivity
content at the archimedean place.
\end{theorem}

\begin{remark}[Deninger's archimedean Iwasawa conjecture]
\label{v4-arith:w8-atp:rem:deninger-archimedean}
The one Iwasawa-adjacent approach not enumerated above is
Deninger's hypothetical archimedean Iwasawa main conjecture
(Deninger 1998 ICM; 2002 \emph{Motives and algebraic cycles}), which
would interpret A-TP as the archimedean limit of a \(p\)-adic Iwasawa
main conjecture. That conjecture is itself RH-equivalent; it does not
reduce the obstruction.
\end{remark}

\subsection{BGS Is Consequent of RH}
\label{v4-arith:w8-atp:negative:BGS}

\begin{theorem}[GUE statistics do not imply A-TP]
\label{v4-arith:w8-atp:thm:BGS-consequent}
\ClaimStatusProvedHere. A Bohigas--Giannoni--Schmit or GUE statement
for ordinates \(\{\gamma_n\}\) does not imply A-TP. Such a statement
concerns limiting correlations of a real sequence. A-TP is positivity
of the full Weil quadratic form on \(\mathrm{PW}(\mathbb R)\), with
the zero side evaluated at
\[
z_\rho=-i(\rho-\tfrac12).
\]
Replacing \(z_\rho\) by a real ordinate \(\gamma_\rho\) is precisely
the zero-line input. Pair-correlation estimates do not supply that
replacement.
\end{theorem}

\begin{remark}[implication direction in \Cref{v4-arith:tfe:thm:stat-upside}]
\label{v4-arith:w8-atp:rem:statistical-upside-unchanged}
The statistical-upside theorem
\Cref{v4-arith:tfe:thm:stat-upside} remains correct: granted the
arithmetic-quantum-chaos hypotheses (C1)--(C5) of
\Cref{v4-arith:tfe:thm:stat-upside}, BGS applied to the 3D ACS bulk
Hamiltonian is a statistical consequence inside that dynamical model.
The direction is from the dynamical hypothesis to statistics, not from
statistics to Weil positivity.
\Cref{v4-arith:w8-atp:thm:BGS-consequent} rules out the latter use.
\end{remark}

\section{The Weil Cone}
\label{v4-arith:w8-atp:residual}

The five source-family boundary of
\Cref{v4-arith:w8-atp:thm:convergence} fixes the exact object which
remains.

\begin{proposition}[five source families do not contract the Weil cone]
\label{v4-arith:w8-atp:prop:consolidated}
\ClaimStatusProvedHereConditional \emph{(on
\Cref{v4-arith:w8-atp:thm:convergence})}. None of the five source
families replaces A-TP by a smaller unconditional theorem. Each
requires the full Weil positivity statement, or a sign condition whose
strength is at least the corresponding A-TP subproblem:
\begin{itemize}
\item Arthur--Selberg retains the archimedean scattering sign.
\item Geometric Langlands retains the arithmetic spectral-action
comparison and the zero-divisor input.
\item Iwasawa retains the absence of an archimedean positivity theorem.
\item Arithmetic quantum chaos retains the zero-line passage from
complex zeros to real ordinates.
\item Arakelov retains Bott--Chern and vertical residual signs.
\end{itemize}
\end{proposition}

\section{Position within the A-TP Obstruction Sequence}
\label{v4-arith:w8-atp:relations}

Chapters~\ref{v4-arith:w5-a:chapter}--\ref{v4-arith:w7-j:ATP-adelic-scattering}
supply ten obstructions on A-TP from analytic directions: Riemann-Siegel,
Nyman-Beurling, de Branges Hilbert-Braun, Gaussian, bounded direct,
Mellin-Barnes, Hadamard, Gram, Guinand-Weil, adelic scattering. They
give sub-class \(W_{\infty}\)-positivity, conditional phase models, or
RH-equivalent reformulations. The
five algebraic traditions below add five further tests from algebraic,
categorical, and Diophantine directions. Two produce new diagnostics
(\Cref{v4-arith:w8-atp:thm:Quant-A-TP-RS},
\Cref{v4-arith:w8-atp:thm:Arakelov-Heegner}); two sharpen the
obstruction by ruling out spurious directions
(\Cref{v4-arith:w8-atp:thm:iwasawa-auxiliary-structure},
\Cref{v4-arith:w8-atp:thm:BGS-consequent}). Across all fifteen
obstructions, the common object is the full Weil positivity problem
\(\mathsf{R}_{\mathrm{A\text{-}TP}}\).

The synthesis of Chapter~\ref{v4-arith:w5-synth:chapter}
(\Cref{v4-arith:w5-synth:thm:RH-equiv-ATP}) establishes
RH \(\Leftrightarrow\) A-TP via the pillar reductions collected there.
This chapter does not further reduce RH below A-TP; it records why the
five source families do not bypass the full Weil cone.

The arithmetic-gravity axis of
Chapters~\ref{v4-arith:3d-acs-E1bar}--\ref{v4-arith:tfe-gutzwiller}
provides two further conditional languages for the same Weil-positive
distribution:
chiral-homology positivity and prismatic Frobenius at the archimedean
place. \Cref{v4-arith:w8-atp:thm:BGS-consequent} constrains the
arithmetic-quantum-chaos statistical-upside theorem
\Cref{v4-arith:tfe:thm:stat-upside}: the implication runs from
quantum chaos to GUE statistics, not from GUE statistics to positivity.

\section{Boundary}
\label{v4-arith:w8-atp:summary}

\begin{enumerate}
\item \textbf{Five source-family boundary
(\Cref{v4-arith:w8-atp:thm:convergence}).} Arthur--Selberg,
geometric Langlands, Iwasawa--Kato, arithmetic quantum chaos, and
Arakelov--Faltings--Hriljac each reaches the Weil positivity problem
or a conditional subproblem. None supplies an unconditional replacement
for A-TP.

\item \textbf{Two conditional diagnostics.}
\Cref{v4-arith:w8-atp:thm:Quant-A-TP-RS}: the Rudnick-Sarnak window is
not an unconditional A-TP theorem.
\Cref{v4-arith:w8-atp:thm:Arakelov-Heegner}: the Heegner-adapted class
reduces Weil positivity to arithmetic Hodge negativity together with
the residual sign \(R_{F,K}(f)\ge0\). The two diagnostics are disjoint:
one is parameterized by pair-correlation support, the other by the
central derivative of an \(L\)-function.

\item \textbf{Two negative results.}
\Cref{v4-arith:w8-atp:thm:iwasawa-auxiliary-structure}: the Iwasawa-Kato
list controls the functional-equation symmetry at finite primes
but carries no positivity information at the archimedean place.
\Cref{v4-arith:w8-atp:thm:BGS-consequent}: GUE statistics for a real
ordinate sequence do not imply Weil positivity and do not put the
zeros on the critical line.

\item \textbf{Uncontracted residual
(\Cref{v4-arith:w8-atp:prop:consolidated}).}
The five source families do not shrink the full Weil cone. The
remaining object is \(\mathsf{R}_{\mathrm{A\text{-}TP}}\), namely
\(W(f*\widetilde f)\geq0\) for all Paley--Wiener \(f\).
\end{enumerate}

The Hilbert--P\'{o}lya obstruction, reformulated as Weil positivity by
Weil 1952, remains the full A-TP residual. A coordinate shadow which
forgets the polar terms, the finite-prime von Mangoldt distribution,
or the full Paley--Wiener cone is not equivalent to RH.
